\section{Propagation and convergence of the rescaled flow}\label{sec:flow}

The localized inequality of \Cref{sec:localized} controls a single complete graph.  We now improve such graphs, propagate them for a fixed time, and sum the resulting motion with radii chosen from the preceding energy drops.

\subsection{Improvement of a graphical slice}

\begin{lemma}[Improvement on a smaller ball]\label{flow:improvement}
Fix a spatial loss $\zeta\in(0,1)$.  Suppose the hypotheses of
Proposition~\ref{flow:localized} hold to radius $R$, and
\begin{equation}\label{flow:small-speed}
 d\le C_0e^{-R^2/8}.
\end{equation}
Suppose also that, on each fixed interior fraction of the graphical
ball, the actual second fundamental form and its derivatives through
a sufficiently large fixed order are bounded by a fixed polynomial
in $R$.  Then the physical height and first derivatives over the
same $S_b$ tend uniformly to zero on $B_{(1-\zeta)R}$ as
$R\longrightarrow\infty$.  The compact multipliers of $S_b$ tend
to zero too.  All statements are uniform in the fixed parameter
neighborhood and in the polynomial geometric bounds.
\end{lemma}

\begin{proof}
Choose a fixed collar containing the entrances to all ends and lying
outside $K_0$, where $S_b$ is an exact shrinker.  Include this collar
and the observation core in a compact reference region $K$.  Since
the conclusion concerns $R\to\infty$, increase the lower radius
threshold to $R_{\eta,K}$ in Proposition~\ref{flow:localized}; its
cutoff $\chi_R$ is then one on both regions.  The compact estimate
and \eqref{flow:base-multiplier} give exponentially small compact error
and multiplier.  Interpolation with the assumed derivatives makes the
height and slope exponentially small on the entrance collar.

We first obtain pointwise control of the forcing on a curtailed
ball.  Fix $\epsilon>0$, and let $q$ have radius at most
$(1-\epsilon)R$.  On a uniform small geometric patch of radius
$\rho$ around $q$, contained in $B_{(1-\epsilon/2)R}$ for large
$R$, \eqref{flow:small-speed} implies
\begin{equation}\label{flow:forcing-L2}
 N_q:=\|\Phi\|_{L^2(\text{patch})}
 \le C\exp\left(-\frac{1-(1-\epsilon/2)^2}{8}R^2\right).
\end{equation}
A bound for $A$ and $\nabla A$ gives
$\|\nabla\Phi\|_\infty\le L_R\le C(1+R)^k$: differentiating
$X^\perp$ introduces at most one position factor.  Put
$A_q=|\Phi(q)|$.  A disk of radius comparable to
$\min\{\rho,A_q/(2L_R)\}$ has $|\Phi|\ge A_q/2$.
The uniform two-dimensional area lower bound on these graphical
patches gives
\begin{equation}\label{flow:forcing-pointwise}
 A_q\le C\bigl(L_R^{1/2}N_q^{1/2}+\rho^{-1}N_q\bigr).
\end{equation}
When $L_R=0$ only the second term is needed.  Polynomial factors
can be absorbed in a smaller positive quadratic exponent.  Thus
$|\Phi|\le Ce^{-cR^2}$ on the curtailed ball.  The exponent
$1/8$ in \eqref{flow:forcing-L2} is the Gaussian norm exponent;
replacing it by the energy exponent $1/4$ would be incorrect.

Let $u$ be the actual normal height over an exterior portion of
$S_b$.  Its equation, divided by the normal component of the
graphical displacement, has the form
\begin{equation}\label{flow:exterior-scalar}
 \begin{gathered}
 \mathcal L_u u=f,\qquad
 \mathcal L_u=a^{ij}\nabla_{ij}
       +\bigl(-\tfrac12Y^\top+\mathfrak b\bigr)\cdot\nabla+c,\\
 |\mathfrak b|\le o(1)(1+|Y|)+C,\qquad
 |c|\le C,\qquad |f|\le C|\Phi|.
 \end{gathered}
\end{equation}
The second-order part is uniformly elliptic and close to the base
metric.  The bounded zeroth-order coefficient deserves verification.
With the convention $D\nu=-A$, the unnormalized graph normal is
$\nu-(I-uA)^{-1}\nabla u$.  Therefore the support-function term,
after this division, is exactly
\begin{equation}\label{flow:support-transport}
 \tfrac12\left(Y\cdot\nu+u
      -Y^\top\cdot(I-uA)^{-1}\nabla u\right).
\end{equation}
Keep its entire transport expression in the first-order part.
In particular, do not differentiate its coefficient in $u$ and
put $|Y|\,|\nabla u|$ into a purported bounded potential.
In the curvature expression, keep the actual coefficients on
$\nabla^2u$ and expand the remaining smooth terms about $u=0$.
Their coefficients depend on the bounded base geometry and the
small first graph variables, and are bounded.  The base equation
vanishes on this annulus.  This proves
\eqref{flow:exterior-scalar} with no compact multiplier forcing.

Write $r=|Y|$.  Uniformly on the far ends, $|Y^\perp|\le C$,
so $|\nabla r|^2=1+O(r^{-2})$ and $\nabla^2r=O(r^{-1})$.
After choosing the weak tolerance small, a sufficiently large
fixed $p$ and a sufficiently small fixed $\alpha_0>0$ give
\begin{equation}\label{flow:supersolutions}
 \mathcal L_u r^p\le-c_1r^p,
 \qquad
 \mathcal L_u e^{\alpha_0r^2}
                   \le-c_2r^2e^{\alpha_0r^2}
\end{equation}
past a fixed entrance.  For the first inequality the leading
coefficient is $-p/2+c+o(p)+O(p^2/r^2)$; for the second it is
$(4\alpha_0^2(1+o(1))-\alpha_0+o(\alpha_0))r^2+O(1)$.
Take, for example, $\alpha_0<1/4$ with a strict margin, then
decrease the graph tolerance and enlarge the entrance.
The comparison principle applies despite the possible positivity
of $c$: dividing a tested function by the positive supersolution
$r^p$ gives a strictly negative zeroth-order coefficient.

On an annulus whose outer radius is $R_1=(1-\epsilon)R$, inner
data and $f$ are exponentially small, while outer data are
bounded by the weak height tolerance.  The function
\[
 V(Y)=C e^{-cR^2}r^p
       +C\varepsilon_{\rm weak}
                     e^{\alpha_0(r^2-R_1^2)}
\]
can, by increasing its first constant, be arranged to dominate
$|u|$ on the two boundaries and to satisfy
$\mathcal L_uV\le-|f|$.  Applying comparison to $V-u$ and
$V+u$ gives $|u|\le V$.  A further fixed fractional curtailment
makes the second summand exponentially small.  The first stays
exponentially small after multiplication by its fixed polynomial.
Interpolation with the available derivative bounds gives the
same conclusion for $\nabla u$.  Allocate the two curtailments
and the small collar margins inside the prescribed $\zeta$.
\end{proof}

\subsection{Finite propagation of the complete graph}

We record carefully the local result used for propagation.  In a
splitting $\mathbb R^{n+k}=\mathbb R^n\times\mathbb R^k$, write
$C_r(x)$ for the cylinder with both horizontal and vertical radii
$r$, centered at $x$.  The graphical pseudolocality theorem of
Ilmanen--Neves--Schulze \cite[Theorem~1.5]{INS2019} states:
for a smooth embedded mean curvature flow with area ratios bounded
by $D$, and every $\upsilon>0$, there are
$\varepsilon,\delta>0$, depending only on $n,k,D,\upsilon$, such
that if $x\in M_0$ and the \emph{whole} intersection
$M_0\cap C_r(x)$ is the graph over the horizontal radius-$r$
ball of a function with Lipschitz constant less than
$\varepsilon$, then $M_t\cap C_{\delta r}(x)$ is a graph on the
full horizontal ball, with Lipschitz constant less than
$\upsilon$ and height at most $\upsilon\delta r$, for
$0\le t<\delta^2r^2$ while the flow exists.  The statement at
general $r$ follows by parabolic scaling.  We use it only for
$n=2,k=1$ and smooth closed flows.

The positive-time input is the graphical interior estimate of
Ecker--Huisken.  For a smooth hypersurface flow remaining graphical
over a fixed horizontal ball $B_a$ throughout $[0,t]$, with compact
graphical closure and $|Df|\le L$, it gives, for $0<\vartheta<1$,
\begin{equation}\label{flow:graphical-smoothing}
 \sup_{B_{\vartheta a}}|\nabla^m A|^2(\cdot,t)
 \le C(m,\vartheta,L)(a^{-2}+t^{-1})^{m+1},\qquad t>0.
\end{equation}
Here $m=0$ follows from
\cite[Corollary~3.2(ii)]{EckerHuisken1991}, and higher orders from
\cite[Theorem~3.4 and Corollary~3.5(ii)]{EckerHuisken1991}.
The constant uses the gradient bound throughout the graphical
spacetime region and no initial curvature norm.  INS supplies that
region and its gradient bound.  We restrict to a smaller horizontal
ball whose graphical closure is compact and stays away from the
vertical cylinder boundary; smoothness and closedness of the actual
flow ensure this local compactness.  Thus
\eqref{flow:graphical-smoothing} applies after each restart with
$t$ measured from that restart.  In particular, for $t\le c a^2$,
$|A|\le Ct^{-1/2}$, and on $t\ge t_*>0$ it gives every fixed
interior derivative bound used below.  The graph equation implies
$|f_t|\le\sqrt{1+L^2}|H|$, so integration gives a uniform height
modulus back to time zero.  Uniform $C^1$ smallness through zero is
proved separately below.  The initial-curvature propagation in
\cite[Remark~1.6(i)]{INS2019} is not needed to supply an initial
curvature bound.  The finite comparison lemma follows by combining
these estimates with overlapping graphical covers and avoidance.

\begin{lemma}[Uniform finite comparison]\label{flow:propagation}
Fix the closed small parameter neighborhood above, an area-ratio
bound, $0<U<1$, a spatial loss $0<\mu<1$, and a target weak $C^1$
normal graph tolerance $\varepsilon_1>0$.  There are
$R_0<\infty$ and $\varepsilon_0,\ell_0>0$ such that the
following holds.  If a smooth closed mean curvature flow $N_t$
has the prescribed area ratios, $N_0$ is a complete normal graph
of norm at most $\varepsilon_0$ over $S_b$ to radius $R\ge R_0$,
and $|\lambda(b,0)|\le\ell_0$, then for $0\le t\le U$
it is a complete graph of norm at most $\varepsilon_1$ over
\[
 S_b(t)=\sqrt{1-t}\,S_b
\]
in $B_{(1-\mu)R}$, while the smooth flow exists.  On every fixed
positive-time interval $[t_*,U]$, all fixed interior curvature
derivative orders are bounded uniformly.  More precisely, the
graph error has a modulus tending to zero as
\[
 \varepsilon_0+|\lambda(b,0)|+R^{-1}\longrightarrow0,
\]
and after positive time this holds in every fixed derivative
order, on the same ball with an arbitrarily small further margin.
The constants are fixed before $R$ or any discrete time index.
In particular the lemma applies with $U=1-e^{-2}$.
\end{lemma}

\begin{proof}
The geometric bounds used here follow uniformly from the finite
parameter family and its end asymptotics.  Shrinking factors in
$[\sqrt{1-U},1]$ preserve bounded geometry and a positive tubular
radius.  Although the raw homothetic parametrization has large
tangential velocity at a distant point, its \emph{normal} velocity
is bounded uniformly, since $|Y^\perp|$ is bounded on $S_b$.
Thus, after tangential reparametrization, comparison motion costs
only a bounded additive spatial buffer on $[0,U]$.

The tubular bound here includes extrinsic separation.  On the compact
core and its collars, proper embeddedness and small smooth variation
in the closed parameter family give a common tubular neighborhood.
On each cylindrical tail the uniformly decaying graph over its
rotated round cylinder retains a fixed tubular radius.  On a conical
tail, rescaled compact annuli are small smooth graphs over annuli of
an embedded cone; the embedded link and its small variations have
uniformly separated nonadjacent pieces.  Scaling back gives a tubular
radius bounded below by a fixed multiple of the distance from the
origin.  The separated end links and directions exclude close pairs
from distinct tails.  The remaining transition region is compact.
These observations give one $\iota>0$ for the complete family,
preserved up to a fixed factor by the shrinking motion on $[0,U]$.
Thus curvature bounds alone are not being used to exclude a second
reference sheet from an entrance cylinder.


We first describe the compactness class of possible comparisons.
If $b_i$ stays in the closed parameter neighborhood and
$\lambda(b_i,0)\to0$, a subsequence has $b_i\to b_\infty$.
At bounded observation centers the limit $S_{b_\infty}$ is a
globally stationary shrinker, because its only possible defect
was the compact source and its multiplier now vanishes.  At
centers going to infinity down a cylindrical end, translation
gives a round cylinder, whose radius at time $t$ is
$\sqrt{2(1-t)}$.  At centers going to infinity down a conical
end, translation gives a plane, stationary under the limiting
motion.  The smooth end estimates and separation of the finitely
many ends give these limits on every fixed ball, in all fixed
orders; another end cannot enter a bounded translated patch.
The bounded-center models are complete shrinkers with bounded
geometry, and all these models have smooth uniformly controlled
motions on $[0,U]$.  No assertion that $b_i\to0$ has been used.

Fix a small INS output slope $\upsilon<1/100$, obtain its entrance
constants $\varepsilon,\delta\in(0,1)$, and choose $r<\iota/100$ so that
every radius-$4r$ reference patch is a graph of slope less than
$\varepsilon/4$.  Set $s=\delta r$.  Choose $\rho<s/100$, then
choose $h>0$ so that
\begin{equation}\label{flow:patch-time}
 h<s^2/8,\qquad Vh<\rho/8,\qquad 4h<\rho^2/32,
\end{equation}
and reference tangent variation during $h$ is small.  Here $V$
bounds reference normal speed.  All choices precede the small
initial-error threshold, which in particular is less than $\rho/8$.
The number $J\le\lceil U/h\rceil+1$ of slabs is fixed.

For $y$ on the reference at the start of a slab, center the entrance
cylinder at the actual point $x=y+u(y)\nu(y)$, with horizontal plane
parallel to $T_yS_b$.  On the larger reference patch, horizontal
projection of the initial normal graph is a small $C^1$ perturbation
of the identity.  Its restriction to a larger horizontal disk is
injective and, by its boundary margin, covers the entire radius-$r$
disk about $x$.  Its height stays inside the vertical cylinder.
Uniform reach excludes a different reference patch; the complete
initial graph and the outer data buffer exclude any other actual
portion.  Thus the \emph{whole} entrance intersection satisfies INS.
Apply it at every such actual center and retain half of each
radius-$s$ output cylinder.

Reserve a fixed additive data buffer, including an extra collar
containing these retained cylinders.  If $p$ in this buffered region
satisfies $\operatorname{dist}(p,S_b(t))\ge\rho$ at $t\le h$,
then its distance from the initial reference is at least
$\rho-Vh$.  Its ball $B_{\rho/2}(p)$ is initially disjoint from
the entire actual surface.  Avoidance with a shrinking sphere,
whose squared radius is $\rho^2/4-4t$, excludes $p$ from the
actual flow.  Every actual point is therefore in the fresh
$\rho$-tube.  If $q$ is its nearest reference point, follow the
reference normal parametrization back to $y$, and choose $x$ as
above.  Then
\begin{equation}\label{flow:refined-membership}
 |p-x|\le |p-q|+|q-y|+|y-x|
          <\rho+Vh+\rho/8<s/4.
\end{equation}
Every actual point consequently lies in an interior retained output.

Normal projection has exactly one preimage on each reference fiber.
To prove existence, fix $q$ and its cylinder as above.  Write its
whole output graph as $z=f(v)$ on the full horizontal disk $B_s$.
Since $|q-x|\le Vh+\rho/8<\rho/4$, the disk used next is
strictly inside $B_s$.
On the disk $B_{s/2}(q_H)$, the horizontal coordinate $G(v)$ of the
normal projection of $(v,f(v))$ differs from $v$ by at most $\rho$,
by confinement and reach.  Its boundary therefore avoids $q_H$,
and the homotopy $v+\sigma(G(v)-v)$ has degree one about $q_H$.
The projected image contains $q$, proving existence.  All the disks
and their graphical lifts lie in the buffered region.  For uniqueness,
two points on that fiber are in the same retained cylinder by
\eqref{flow:refined-membership}.  If their normal parameters differ
by $a$, their vertical separation is at least
$\sqrt{1-\omega^2}|a|$ and their horizontal separation is at most
$\omega|a|$, where $\omega\ll1$ is the reference normal's tilt.
The actual graph instead bounds the vertical separation by
$\upsilon\omega|a|$, forcing $a=0$.  The same small slopes make
normal projection a local diffeomorphism.  Full-disk coverage thus
gives existence, and the slope bound gives uniqueness.  The local
graphs agree on overlaps and cover artificial patch edges.  Notice
that no inequality $\upsilon s<\rho$ was assumed: avoidance
supplies the sharper normal height independently of the INS height
bound.

We verify the vanishing error modulus also at elapsed times tending
to zero.  Suppose a sequence with initial graph error, compact
multiplier, and inverse data radius tending to zero has graph error
at least $\eta>0$ at times $t_i$.  Let $K$ bound reference curvature
and $C_0$ the normal-chart conversion constants.  For the case $t_i\to0$, make
the following \emph{fresh} choices for this target $\eta$.  First
choose an arbitrarily small INS output slope $\upsilon$, then its
constants $\varepsilon,\delta$.  Next choose $r\le1$ below the
uniform reach scale so that the reference entrance slope is below
$\varepsilon/4$ and its tangent variation $Kr$ is as small as
required.  Put $s=\delta r$, choose
\[
 \rho<\min\{\iota/100,s/100,\upsilon s/100,\eta/(100C_0)\},
\]
and only then choose $h_\eta>0$ satisfying
\eqref{flow:patch-time} and making the uniform reference tangent
modulus $\omega_{\rm ref}(h_\eta)$ small.  Arrange
\[
 C_0\bigl(\upsilon+Kr+\omega_{\rm ref}(h_\eta)\bigr)<\eta/2.
\]
Only after these choices take $i$ large enough that the initial
error meets all entrance and $\rho/8$ margins, the additive buffer
lies inside the data region, and $t_i<h_\eta$.

Repeat the shrinking-sphere confinement at this new $\rho$, rather
than using the earlier, wider tube.  Equation
\eqref{flow:refined-membership} now places every actual point in a
refined output cylinder, including the purported point of failure.
Its height over the reference is less than $\rho$, and its slope
in the normal chart is bounded by
$C_0(\upsilon+Kr+\omega_{\rm ref}(h_\eta))$.  Both are smaller
than the alleged error threshold.  This excludes $t_i\to0$.
No relation between $\delta$ and $\upsilon$ was needed, and no
patch radius was chosen as a function of the failing time.

If instead $\liminf t_i>0$, recenter at the points of failure.
The data domains exhaust, and the complete graphs, area bound, and
\eqref{flow:graphical-smoothing} give a smooth limiting motion for
positive time.  The integrable curvature-speed bound gives its
initial height continuously.  The limiting comparison is an exact
smooth mean curvature flow: the base multipliers vanish in the
bounded-center limit, and the end limits are the shrinking cylinder
or stationary plane.  We identify these two motions as follows.

Use a normal-velocity parametrization $F(y,t)$ of the complete
comparison, and write the limiting flow as $F+u\nu$.  Its equation
in the reference metric has the form
\[
 u_t=a^{ij}(y,t,u,\nabla u)\nabla_{ij}u
              +B(y,t,u,\nabla u),\qquad B(y,t,0,0)=0.
\]
The principal coefficients are uniformly elliptic and bounded.
The lower function $B$ has bounded first derivatives in its last
arguments, since it involves only bounded reference geometry and
motion, the small height, and its first derivative.  Retain the
\emph{actual} $a^{ij}$ and linearize only $B$:
\[
 B=cu+b^i\nabla_i u,\qquad
 (c,b)=\int_0^1(B_u,B_{\nabla u})(y,t,\sigma u,\sigma\nabla u)
                                                       \,d\sigma.
\]
This gives bounded lower coefficients without multiplying an unknown
Hessian by derivatives of $a^{ij}$, even near time zero.  The exact
comparison equation is used in $B(y,t,0,0)=0$.
For $\psi=1+|F|^2$, bounded normal speed and reference geometry give
$|\partial_t\psi|+|\nabla\psi|+|\nabla^2\psi|\le C\psi$.
Hence $\epsilon e^{Ct}\psi$ is a supersolution for both signs of
the bounded height $u$, for one sufficiently large $C$.
Properness supplies compact exhaustions, and continuous zero initial
height supplies their initial data.  The maximum principle, followed
by exhaustion and $\epsilon\downarrow0$, gives $u=0$.
Normal parametrization has removed the raw homothetic tangential
velocity; the needed change of parameters exists on the finite
horizon because that tangential field has at most linear growth.
This contradicts positive-time failure and proves the modulus on
the whole first slab.

On the latter half of a slab, the interior estimates and this
identification give smooth closeness to the comparison.  Its complete
normal graph is the \emph{union} of the overlapping outputs.
A fresh radius-$4r$ reference patch in the smaller data region is
covered by finitely many of their interiors, with a number depending
only on the fixed ratio $r/s$.  Thus its fresh radius-$r$ entrance
cylinder contains the whole actual intersection, with the required
small entrance slope.  Center it at the corresponding actual graph
point.  A single old radius-$s$ output is never asked to contain the
new larger entrance cylinder.

Keep the original coarse $r,h$ and weak tolerances for these
restarts.  The refined choices above serve only to test the modulus
at a slab entrance.  Finite induction over
$J\le\lceil U/h\rceil+1$ selects initial and multiplier tolerances
so that each restart is admissible and its error tends to zero.
Every slab consumes a fixed additive buffer $A$ for enlarged
patches, motion, and the confinement collar.  Thus the total loss is
at most $JA$, independent of the covering cardinality, $R$, and
the later discrete index.  Taking $JA<\mu R$ proves the asserted
complete graph on $B_{(1-\mu)R}$.  Positive-time estimates give all
fixed derivative bounds and the corresponding smooth error modulus
away from elapsed time zero.
\end{proof}

\subsection{Delayed radii and summability}

The remaining task is to use the localized inequality for arbitrarily
long times.  The graphical radius cannot be chosen independently at
each slice, since its next value must be reached by propagation from
the previous slice.  We therefore choose each radius from all preceding
energy drops.  This gives both the propagation recurrence and a
summable bound for the accumulated Gaussian errors.

We shift only the origin of the rescaled clock so that the initial
time $0$ is a sufficiently late member of the chosen tangent
sequence.  On any prescribed fixed spatial ball and bounded
initial time window the flow is then as smoothly close to $S$
as desired, and $\mathcal E(0)$ is as small as desired.  This follows from
the smooth multiplicity-one convergence in
Proposition~\ref{geo:tangent}, where the hypotheses of
multiplicity-one local regularity are verified.  The initial
window will include two rescaled time units.

Fix a sufficiently large observation ball and a larger fixed
collar containing the end entrances used in Lemma~\ref{flow:improvement}.
Let $K$ be the corresponding compact reference region.  Bootstrap the
condition that the actual flow is a
complete graph in a small weak $C^1$ neighborhood of $S$ on
this ball and collar.  This makes the observations and the
coordinates $(b,x)$ well defined.  All radius statements in
the following induction are made while this core bootstrap
holds; its closure will be proved below.

Choose $\eta,c$ with $0<\eta<c<\theta$, where $\theta$ is fixed
by Proposition~\ref{flow:localized}.  Choose a small spatial
loss $\zeta>0$, and then choose $Q$ with
\begin{equation}\label{flow:Q-choice}
 1<Q<(1-2\zeta)e^{0.4}.
\end{equation}
All graphical tolerances, derivative orders, and the finite
propagation horizon $U=1-e^{-2}$ are now fixed.  Finally choose
the minimum radius at least $R_{\eta,K}$, so that the observation core
and entrance collar lie in the region where $\chi_R=1$, and otherwise
sufficiently large; then choose $\mathcal E(0)$ sufficiently
small.  This ordering permits all improvements and propagation
conclusions to land strictly inside their weak tolerances.

For each unit interval wholly available in the bootstrap put
\begin{equation}\label{flow:radii}
 \begin{gathered}
 a_j=(\mathcal E(j)-\mathcal E(j+1))^{1/2},\quad
 r_j=\sqrt{8\log(1/a_j)},\\
 R_j=\min\left\{Q^jR_{\rm in},
                 \min_{0\le i\le j}Q^{j-i}r_i\right\}.
 \end{gathered}
\end{equation}
Here $r_j=+\infty$ if $a_j=0$; we choose $\mathcal E(0)<1$ so that
all other logarithms are positive.  Select a time
\begin{equation}\label{flow:selected-time}
 \tau_j\in[j+0.2,j+0.4],\qquad
 d(\tau_j)\le\sqrt5\,a_j
                     \le\sqrt5\,e^{-R_j^2/8},
\end{equation}
by averaging $d^2$ on this subinterval.  In the zero case the
speed vanishes there and any interior time is suitable.  The
radius definition gives
\begin{equation}\label{flow:radius-recursion}
 R_{j+1}=\min\{QR_j,r_{j+1}\},\qquad
 R_j\ge\min\{R_{\rm in},\sqrt{8\log(1/\sqrt{\mathcal E(0)})}\}.
\end{equation}

\begin{lemma}[Graphical radius induction]\label{flow:radius-induction}
On every complete bootstrap interval, the hypotheses of
Proposition~\ref{flow:localized} and of
Lemma~\ref{flow:improvement} hold at $\tau_j$ to radius $R_j$,
with the harmless outer margins specified above.
\end{lemma}

\begin{proof}
The chosen initial tangent window supplies the assertion for
$j=0$, including the derivative bounds.  Suppose it holds at
$\tau_j$.  The localized inequality, the multiplier estimate,
and \eqref{flow:selected-time} make the core height and compact
multiplier small as the minimum radius becomes large.
Lemma~\ref{flow:improvement} gives as small a height and slope
as desired over $S_{b_j}$ to radius $(1-\zeta)R_j$, where
$b_j=b(\tau_j)$.  In particular this smallness is much stronger
than the weak tolerance being propagated.

Normalize the unrescaled time at this slice to $-1$.
For a rescaled duration $\Delta$ the ordinary elapsed time is
$1-e^{-\Delta}$, and dilation back to the rescaled coordinates
has factor $e^{\Delta/2}$.  Consecutive selected times satisfy
$0.8\le\tau_{j+1}-\tau_j\le1.2$.  Apply
Lemma~\ref{flow:propagation}, allocating its additional spatial
loss within another $\zeta R_j$.  In the rescaled coordinates
it yields a complete weak graph relative to the old $S_{b_j}$
to radius at least
\begin{equation}\label{flow:radius-gain}
 (1-2\zeta)e^{(\tau_{j+1}-\tau_j)/2}R_j,
\end{equation}
and throughout the intervening times on the corresponding
interior regions.  Its positive-time smoothing supplies the
derivative bounds at the next selected slice before any new
elliptic improvement is invoked.  If an arbitrarily small
additional outer margin is needed for those bounds, reserve
it in the strict inequality \eqref{flow:Q-choice}.

It remains to change to the new base $S_{b_{j+1}}$.  While the
fixed core is a weak graph over $S$, the graph equation and
Cauchy--Schwarz on the compact support of the observations give
\begin{equation}\label{flow:observation-motion}
 \left|\frac{d}{d\tau}(b(\tau),x(\tau))\right|
       \le C d(\tau),\qquad
 |b_{j+1}-b_j|\le C(a_j+a_{j+1}).
\end{equation}
To see the first inequality, the fixed-core normal graph height
$v$ satisfies $\partial_\tau v=\Phi\cdot\nu_M/(\nu_S\cdot\nu_M)$;
the denominator stays bounded away from zero.  Each observation
is a fixed compact integral of $v$, and the inverse coordinate
Jacobian is bounded.  The second inequality follows by
integrating on the portions of $[j,j+1]$ and $[j+1,j+2]$ between
the two slices.

Smooth dependence of the rotations, conical links, and decaying
remainders gives a chart-change cost in physical height and
slope at radius $R$ bounded by
$C(1+R)|b_{j+1}-b_j|$.  Both adjacent cutoffs in
\eqref{flow:radii} are necessary here.  Specifically
\[
 R_{j+1}a_j\le Q a_j\sqrt{8\log(1/a_j)},\qquad
 R_{j+1}a_{j+1}\le a_{j+1}\sqrt{8\log(1/a_{j+1})},
\]
with the products interpreted as zero at $a_i=0$.
These are uniformly small when $\sup a_i\le\sqrt{\mathcal E(0)}$ is
small.  Thus the new-chart graph still lies inside its weak
tolerance.  Equations \eqref{flow:Q-choice},
\eqref{flow:radius-recursion}, and \eqref{flow:radius-gain}
place $B_{R_{j+1}}$, with its required margins, in that graph.
The next localized inequality and improvement can now be used.
This proves the induction in the order: previous improvement,
parabolic propagation and smoothing, chart change, then new
elliptic improvement.
\end{proof}

The lemma and monotonicity, since $\tau_j<j+1$, imply
\begin{equation}\label{flow:discrete-gradient}
 \mathcal E(j+1)^{1-\theta}
 \le C\left(a_j+e^{-(1-\eta)R_j^2/8}\right).
\end{equation}
The following summation keeps the gradient-scale tail in this
exact form.

\begin{lemma}[Uniform finite-prefix length estimate]
\label{flow:finite-length}
Put
\[
 p=\frac{\theta-c}{1-c}\in(0,1/2),\qquad
 \kappa=\frac{1-\eta}{1-c}>1,\qquad q=Q^2,
 \qquad h=\kappa/8.
\]
For $\mathcal E(0)$ sufficiently small, every available finite prefix
satisfies
\begin{equation}\label{flow:length-bound}
 \sum_{j=0}^{n}a_j
 \le 2\bigl(C\mathcal E(0)^p+B(R_{\rm in})\bigr),\qquad
 B(R)=\frac{e^{-hR^2}}{1-e^{-h(q-1)R^2}}.
\end{equation}
In particular the length on complete unit intervals can be
made arbitrarily small, uniformly in the number of intervals.
\end{lemma}

\begin{proof}
Call $j$ good if
$e^{-(1-\eta)R_j^2/8}\le a_j^{1-c}$, and bad otherwise.
For a good index with $a_j>0$, \eqref{flow:discrete-gradient}
and $a_j\le1$ imply
\[
 \mathcal E(j+1)^{1-p}\le C_1a_j.
\]
Write $u=\mathcal E(j)$ and $v=\mathcal E(j+1)$.  If $v\ge u/2>0$,
concavity and the last inequality give
\[
 u^p-v^p\ge p u^{p-1}(u-v)
          =p u^{p-1}a_j^2\ge c_1a_j.
\]
If $v<u/2$, then, since $u\le1$ and $p<1/2$,
\[
 u^p-v^p\ge(1-2^{-p})u^p
           \ge(1-2^{-p})\sqrt u\ge(1-2^{-p})a_j.
\]
Zero terms cause no difficulty.  Summing the nonnegative energy
differences only over good indices gives
$\sum_{j\text{ good}}a_j\le C\mathcal E(0)^p$ on every prefix.

At a bad index,
\begin{equation}\label{flow:bad-max}
 a_j<e^{-hR_j^2}
 =\max\left\{e^{-hq^jR_{\rm in}^2},
                \max_{i\le j}a_i^{\kappa q^{j-i}}\right\}.
\end{equation}
If $0<a_j<1$, its own term $a_j^\kappa$ is strictly less
than $a_j$ and cannot account for this maximum.  If $a_j=0$
the following bound is immediate.  Thus
\begin{equation}\label{flow:bad-convolution}
 a_j\le e^{-hq^jR_{\rm in}^2}
                  +\sum_{i<j}a_i^{\kappa q^{j-i}}.
\end{equation}
This exclusion of $i=j$ is where $\kappa>1$ is essential.

Let $\delta_0=\sqrt{\mathcal E(0)}<1$ and
\[
 \rho(\delta_0)=\sum_{\ell\ge1}
                   \delta_0^{\kappa q^\ell-1}
 \le\frac{\delta_0^{\kappa q-1}}
              {1-\delta_0^{\kappa(q-1)}}\longrightarrow0.
\]
Since every $a_i\le\delta_0$, summing
\eqref{flow:bad-convolution} on a prefix and interchanging
the finite nonnegative sums bounds its double sum by
$\rho(\delta_0)\sum_{i=0}^n a_i$.  Also
$\sum_{j\ge0}e^{-hq^jR^2}\le B(R)$, using
$q^j\ge1+j(q-1)$.  Choose $\mathcal E(0)$ so small that
$\rho(\delta_0)\le1/2$.  Absorbing proves
\eqref{flow:length-bound}.  Finally
$\int_j^{j+1}d\,d\tau\le a_j$ by
\eqref{flow:energy} and Cauchy--Schwarz.
\end{proof}

\subsection{Closing the bootstrap and fixing every negative-time measure}

\begin{proposition}[Fixed-frame convergence and uniqueness]
\label{flow:uniqueness}
The rescaled flow converges smoothly on every fixed compact set,
at all rescaled times, to $S$ in the original ambient coordinates.
For the smooth closed embedded surface flow and fixed first singular
point of \Cref{thm:main}, every backward tangent flow consequently equals
the multiplicity-one homothetic motion $\sqrt{-s}\,S$ at every
negative time $s$.
\end{proposition}

\begin{proof}
Suppose first that the weak core bootstrap has a finite first
stopping time $L$.  Use only intervals $[i,i+1]$ whose right
endpoints lie strictly before $L$.  If $i$ is the last such
index, then
\begin{equation}\label{flow:last-gap}
 i+1<L\le i+2,
 \qquad \tau_i\ge i+0.2,
 \qquad L-\tau_i\le1.8<2.
\end{equation}
Allowing a completed interval ending at $L$ by continuity only
shortens the gap.  If no such interval is available, the
initial tangent-convergence window supplies the graphs and
derivatives through this prospective stopping time.

For a last available slice, propagate directly from its improved
graph using Lemma~\ref{flow:propagation} over the fixed horizon
$U=1-e^{-2}<1$.  No subsequent selected slice, localized
inequality, or future unit energy drop is needed.  Even a
limiting cylinder retains radius at least $\sqrt2 e^{-1}$ on
this horizon, so all patch constants and buffers were fixed
with a genuine pre-extinction margin.  Choose the minimum
graphical radius so large that this finite additive buffer
leaves the fixed bootstrap ball and its surrounding collar
strictly inside the propagated domain.  This gives graphical
coverage there through $L$.

Uniform derivatives near the last selected time require a preceding
positive delay.  If $i\ge1$, use the propagation from $\tau_{i-1}$ on
\[
 [\tau_i,\min\{L,\tau_i+0.2\}],
\]
where the rescaled delay lies in $[0.8,1.4]$, and use the propagation
from $\tau_i$ on $[\tau_i+0.2,L]$, where it lies in $[0.2,1.8]$.
Both propagations are available on their entire two-unit horizons
from the already proved improved graphs; the earlier induction's use
of part of a horizon does not shorten it.  On every completed selected
interval $[\tau_j,\tau_{j+1}]$ with $j\ge1$, use $\tau_{j-1}$ for
the first $0.2$ time units and $\tau_j$ thereafter.  The initial
two-unit tangent window supplies the first entrance and the cases
$i=0$ or no complete interval.  In these covers the ordinary elapsed
times are bounded below by $1-e^{-0.2}$, and the dilation factors are
uniformly bounded.  Lemma~\ref{flow:propagation} therefore supplies
uniform interior curvature derivatives on the fixed ball and collar.
The weak fixed-core graph and the bounded geometry of $S$ convert
these geometric estimates to uniform derivatives of its normal
height, as needed for interpolation below.  This uses only known
selected slices and no future energy drop.  The complete graphical
coverage also excludes an entering extra sheet or a loss at the
boundary of the observation domain at the stopping time.

The length through completed intervals is bounded by
Lemma~\ref{flow:finite-length}.  The entire remaining segment
from the last selected slice costs at most
\begin{equation}\label{flow:last-length}
 \int_{\tau_i}^L d(\tau)\,d\tau
 \le\sqrt2\left(\int_{\tau_i}^Ld(\tau)^2\,d\tau\right)^{1/2}
 \le\sqrt2 \mathcal E(0)^{1/2}.
\end{equation}
Counting part of a complete interval twice only weakens this
upper bound.  The initial-window case has the same estimate
with its bounded duration.  The graph-speed formula used in
\eqref{flow:observation-motion} therefore bounds the change
of the fixed-core graph in $L^2$ by
\[
 C\bigl(\mathcal E(0)^p+B(R_{\rm in})+\mathcal E(0)^{1/2}\bigr).
\]
The initial graph can be chosen much closer to $S$ than the
bootstrap tolerance.  Interpolation with the uniform interior
derivatives on the larger collar makes the resulting $C^1$
distance strictly smaller than that tolerance through $L$.
The already established coverage and these strict estimates
extend the bootstrap past $L$, a contradiction.  This proves
that the bootstrap and the radius induction hold for all time.

Letting the prefix tend to infinity in
\eqref{flow:length-bound} gives
\begin{equation}\label{flow:finite-total-length}
 \int_0^\infty d(\tau)\,d\tau<\infty,
 \qquad \sum_{j\ge0}a_j<\infty.
\end{equation}
On the fixed observation core the graph functions are Cauchy
in $L^2$, and the uniform derivative bounds make them converge
smoothly on a slightly smaller core.  Their limit is $S$,
because the original tangent sequence still has a subsequence
at arbitrarily large rescaled times and converges there to
$S$.  In particular \eqref{flow:observation-motion} gives
$(b_j,x_j)\to(0,0)$ in the fixed compact observation chart.

The radii actually diverge.  Since $a_i\to0$, their logarithmic
cutoffs $r_i$ tend to infinity.  Given $H<\infty$, choose $I$
such that $r_i\ge H$ for $i\ge I$.  Every recent term
$Q^{j-i}r_i$, $I\le i\le j$, is then at least $H$; each of
the finitely many earlier terms and $Q^jR_{\rm in}$ eventually
exceeds $H$.  Taking the minimum in \eqref{flow:radii} proves
$R_j\to\infty$.

The improved graphs at $\tau_j$ thus have height and slope
tending to zero on every fixed compact subset of $S_{b_j}$.
Their base compact multipliers tend to zero by
\eqref{flow:base-multiplier}.  Apply the error-modulus version
of Lemma~\ref{flow:propagation} on every interval
$[\tau_j,\tau_{j+1}]$.  In the rescaled coordinates its
homothetic comparison is exactly the fixed set $S_{b_j}$,
because the comparison shrinkage is canceled by the rescaling.
Since $b_j\to0$, the actual graphs converge to $S$ at
\emph{every intermediate time}, uniformly on each compact
spatial set.  The higher derivative upgrade always uses a preceding slice with
positive delay.  For $\tau\in[\tau_j,\tau_{j+1}]$, use the slice
$\tau_j$ if $\tau-\tau_j\ge0.2$, giving rescaled delay in
$[0.2,1.2]$.  If $0\le\tau-\tau_j<0.2$, use $\tau_{j-1}$,
giving delay in $[0.8,1.4]$.  Both delays are bounded below by
$0.2$ and above by the fixed horizon $2$, and therefore their
ordinary elapsed times are bounded below by $1-e^{-0.2}>0$
and above by $1-e^{-2}<1$.  The dilation factors are uniformly
bounded.  Since the preceding graphical radii tend to infinity,
every fixed compact set and its interior margin are eventually
contained in these propagated domains.  Equation
\eqref{flow:graphical-smoothing} supplies uniform bounds in every
fixed derivative order there; interpolation with the established
$C^1$ convergence gives smooth convergence, including at each
selected endpoint.  No curvature estimate at zero elapsed time is
used.  Thus convergence holds in the original ambient frame at
every rescaled time.

Restore the original, unshifted rescaled clock.  For every
$s<0$ and every admissible $\lambda$ there is the exact identity
\begin{equation}\label{flow:negative-time-identity}
 M_s^\lambda
 =\sqrt{-s}\,M\bigl(2\log\lambda-\log(-s)\bigr).
\end{equation}
The rescaled times on the right tend to infinity, uniformly
for $s$ in any compact interval $[-B,-b]\subset(-\infty,0)$.
Consequently the full family of actual surfaces on the left
converges locally smoothly and with multiplicity one to
$\sqrt{-s}\,S$.  Smooth complete graphical convergence gives
both local area-measure and varifold convergence for every
individual $s<0$.  The identity uses only the original fixed
center and scalar dilation, so it preserves the axes and all
ambient positions.  This proves the asserted smooth convergence.
To pass from convergence of the actual slices to every slice of a
weak representative, we use the all-endpoint convention and the
weak-limit premise specified below.
\end{proof}

\begin{proof}[Completion of the proof of Proposition~\ref{flow:uniqueness}]
Here is the precise weak-limit premise sufficient for the representative
assertion.  Write
$\mu_s^\lambda=\mathcal H^2\!\llcorner M_s^\lambda$.  A tangent
Radon family $(\mu_s)_{s<0}$ is assumed to satisfy the local hypotheses
and the all-endpoint inequality of Lemma~\ref{geo:representative}, and
along its extracted scales $\lambda_i\to\infty$ to satisfy
\begin{equation}\label{flow:spacetime-weight-limit}
 \int\!\!\int\psi(X,s)\,d\mu_s^{\lambda_i}(X)\,ds
 \longrightarrow
 \int\!\!\int\psi(X,s)\,d\mu_s(X)\,ds
 \quad\bigl(\psi\in C_c(\mathbb R^3\times(-\infty,0))\bigr).
\end{equation}
For each test, the left side is considered once its time support lies
in the interval on which the rescaled flow is defined.
The local mass condition in the lemma makes the right-hand spacetime
weight a locally finite Radon measure.  This is the only weak-limit
identification used here.  Local weak convergence
$\mu_s^{\lambda_i}\rightharpoonup\mu_s$ for almost every $s$, together
with a uniform bound for $\mu_s^{\lambda_i}(B_R)$ on each compact time
interval and for each fixed $R$, also implies
\eqref{flow:spacetime-weight-limit} by dominated convergence.

The direct smooth convergence already proved shows that the left side
of \eqref{flow:spacetime-weight-limit} converges to
$\int\!\!\int\psi\,d\nu_s\,ds$, with
$\nu_s=\mathcal H^2\!\llcorner(\sqrt{-s}\,S)$.  To justify the time
integration, place the support of $\psi$ in $B_R\times J$ with
$J\Subset(-\infty,0)$.  The scale-invariant area-ratio bound gives
$\mu_s^{\lambda_i}(B_R)\le\Lambda R^2$, so dominated convergence
applies to the slice integrals.  Uniqueness of the local weak Radon
limit therefore gives
$\mu_s\,ds=\nu_s\,ds$ as spacetime measures.

Choose a countable family of compactly supported smooth spatial tests
that determines Radon measures, for example a family dense on each
ball in the uniform norm.  Test this spacetime equality by the product
of each spatial test and an arbitrary compactly supported continuous
time test.  Local integrability implies equality of the corresponding
slice integrals almost everywhere for each spatial test.  Intersecting
these countably many full-measure sets and using the determining
property yields $\mu_s=\nu_s$ as Radon measures for almost every $s$.
Lemma~\ref{geo:representative} now gives this equality for every
$s<0$.  The argument applies to every all-endpoint representative
satisfying \eqref{flow:spacetime-weight-limit}, so no choice of its
exceptional slices can change the negative-time measures.  This
proves the representative conclusion under the stated convention.
\end{proof}

\begin{proof}[Proof of Theorem~\ref{thm:main}]
Fix a singular point at the first singular time.  The preceding
construction applies to its tangent section from \Cref{geo:tangent},
and Proposition~\ref{flow:uniqueness} gives the full fixed-frame smooth
convergence and the equality of every negative-time measure under the
stated tangent convention.  The point was arbitrary.
\end{proof}
